\section{Рефакторинг заданий ОГЭ и ЕГЭ}\label{4sect}
В процессе разработки возникает необходимость постоянного совершенствования качества и актуализации исходного кода. В контексте проекта генерации экзаменационных заданий рефакторинг приобретает особую значимость, поскольку от качества кода напрямую зависят корректность сгенерированных задач, возможность их масштабирования и адаптации под изменяющиеся требования экзаменационных спецификаций.
В рамках настоящей главы представлен вклад автора в рефакторинг программного кода проекта <<Час ЕГЭ>>. Рассматриваются конкретные примеры преобразования исходного кода генераторов математических задач, демонстрирующие применение ключевых техник рефакторинга.

\subsection{Вклад автора в рефакторинг на проекте}

\lstinputlisting[caption = refactoring2.js, label={lst:refactoring2}, escapechar=|]{code/refactoring2.js}
\subsubsection*{Примеры генерируемых задач refactoring2.js}
 
\par{Прямая $y=-x+12$ является касательной к графику функции $y=2x^{2}+bx+44$. Найдите $b$, зная, что оно больше $−14$. }{ 15 }{refactoring2.js}
\par{Прямая $y=−56x+95$ является касательной к графику функции $y=-5x^{2}+bx-340$. Найдите $b$, зная, что оно больше $−58$. }{ -126 }{refactoring2.js}
\par{Прямая $y=-2x+11$ является касательной к графику функции $y=2x^{2}+bx+39$. Найдите $b$, зная, что оно больше $−13$. }{ -22 }{refactoring2.js}

\lstinputlisting[caption = refactoring3.js, label={lst:refactoring3}, escapechar=|]{code/refactoring3.js}
\subsubsection*{Примеры генерируемых задач refactoring3.js}
 
\par{Прямая $y=dx+209$ является касательной к графику функции $y=6x{2}+25x+593$. Найдите d. }{ -71 }{refactoring3.js}
\par{Прямая $y=7x+23$ является касательной к графику функции $y=−7x{2}−49x+c$. Найдите c. }{ -89 }{refactoring3.js}
\par{Прямая $y=−33x+66$ является касательной к графику функции $y=−7x{2}−89x+c$. Найдите c. }{ -46 }{refactoring3.js}

\lstinputlisting[caption = refactoring5.js, label={lst:refactoring5}, escapechar=|]{code/refactoring5.js}
\subsubsection*{Примеры генерируемых задач refactoring5.js}
 
\par{В цилиндрическом сосуде уровень раствора достигает $4$ см. На какой высоте будет находиться уровень раствора, если перелить содержимое первого сосуда во второй цилиндрический сосуд, площадь основания которого в $25$ раз меньше площади основания первого? Ответ выразите в сантиметрах. }{ 100 }{refactoring5.js}
\par{В сосуде, имеющем форму правильной треугольной призмы, уровень ртути достигает $36$ см. На какой высоте будет находиться уровень ртути, если перелить содержимое первого сосуда во второй сосуд, имеющий форму правильной треугольной призмы, сторона основания которого в $3$ раза больше стороны основания первого? Ответ выразите в сантиметрах. }{ 4 }{refactoring5.js}
\par{В сосуде, имеющем форму правильной четырёхугольной призмы, уровень уксуса достигает $12$ см. На какой высоте будет находиться уровень уксуса, если перелить содержимое первого сосуда во второй сосуд, имеющий форму правильной четырёхугольной призмы, сторона основания которого в $3$ раза меньше стороны основания первого? Ответ выразите в сантиметрах. }{ 108 }{refactoring5.js}

\lstinputlisting[caption = refactoring6.js, label={lst:refactoring6}, escapechar=|]{code/refactoring6.js}
\subsubsection*{Примеры генерируемых задач refactoring6.js}
 
\par{Два доисторических омнибуса одновременно отправились в $15$-километровый пробег. Первый ехал со скоростью, на $2$ км/ч большей, чем скорость второго, и прибыл к финишу на $2$ часа раньше второго. Найти скорость доисторического омнибуса, пришедшего к финишу вторым. Ответ дайте в км/ч. }{ 3 }{refactoring6.js}
\par{Два "Москвича" одновременно отправились в $48$-километровый пробег. Первый ехал со скоростью, на $6$ км/ч большей, чем скорость второго, и прибыл к финишу на $4$ часа раньше второго. Найти скорость "Москвича", пришедшего к финишу первым. Ответ дайте в км/ч. }{ 12 }{refactoring6.js}
\par{Два "Запорожца" одновременно отправились в $28$-километровый пробег. Первый ехал со скоростью, на $7$ км/ч большей, чем скорость второго, и прибыл к финишу на $2$ часа раньше второго. Найти скорость "Запорожца", пришедшего к финишу вторым. Ответ дайте в км/ч. }{ 7 }{refactoring6.js}

\lstinputlisting[caption = refactoring7.js, label={lst:refactoring7}, escapechar=|]{code/refactoring7.js}
\subsubsection*{Примеры генерируемых задач refactoring7.js}
 
\par{Моторная лодка прошла против течения реки и вернулась в пункт отправления. Известно, что скорость течения равна $3$ км/ч. Насколько меньше времени затратила лодка на обратный путь, если скорость лодки в неподвижной воде составляет $10$ км/ч, а суммарное пройденное лодкой расстояние составляет $364$ км? }{ 12 }{refactoring7.js}
\par{Моторная лодка прошла против течения реки и вернулась в пункт отправления. Чему равно суммарное пройденное лодкой расстояние, выраженное в км, если скорость лодки в неподвижной воде составляет $12$ км/ч, а лодка затратила на обратный путь на $25$ часов меньше? Скорость течения составляет $8$ км/ч. }{ 250 }{refactoring7.js}
\par{Моторная лодка прошла против течения реки и вернулась в пункт отправления. Чему равна скорость лодки в неподвижной воде, выраженная в км/ч, если суммарное пройденное лодкой расстояние равно $30$ км, лодка затратила на обратный путь на $14$ часов меньше? Скорость течения составляет $7$ км/ч. }{ 8 }{refactoring7.js}

\lstinputlisting[caption = refactoring8.js, label={lst:refactoring8}, escapechar=|]{code/refactoring8.js}
\subsubsection*{Примеры генерируемых задач refactoring8.js}
 
\par{Перед началом первого тура чемпионата по вольной борьбе участниц разбивают на игровые пары случайным образом с помощью жребия. Всего в чемпионате участвует $326$ спортсменок, среди которых $79$ участниц из Венесуэлы, в том числе Мария. Найдите вероятность того, что в первом туре Мария будет играть с какой-либо спортсменкой не из Венесуэлы. }{ 0.76 }{refactoring8.js}
\par{Перед началом первого тура чемпионата по настольному теннису участниц разбивают на игровые пары случайным образом с помощью жребия. Всего в чемпионате участвует $566$ спортсменок, среди которых $114$ участниц из Англии, в том числе Екатерина. Найдите вероятность того, что в первом туре Екатерина будет играть с какой-либо спортсменкой из Англии. }{ 0.2 }{refactoring8.js}
\par{Перед началом первого тура чемпионата по бадминтону участниц разбивают на игровые пары случайным образом с помощью жребия. Всего в чемпионате участвует $726$ спортсменок, среди которых $465$ участниц из Австрии, в том числе Елена. Найдите вероятность того, что в первом туре Елена будет играть с какой-либо спортсменкой из Австрии. }{ 0.64 }{refactoring8.js}

\lstinputlisting[caption = refactoring15.js, label={lst:refactoring15}, escapechar=|]{code/refactoring15.js}
\subsubsection*{Примеры генерируемых задач refactoring15.js}
 
\par{Между остановками Ш. и П. $7$ километров прямой трассы. Наталия выехала на велосипеде из леса между станциями в $700$ метрах от Ш. и увидела, что к Ш. в направлении к П. с постоянной скоростью подъезжает автобус, на который Наталии нужно успеть. Она заметила, что если она сейчас поедет в сторону Ш., она окажется там одновременно с автобусом. Но и если она поедет в сторону П., она также окажется там одновременно с автобусом, который успеет преодолеть весь участок от Ш. до П., не останавливаясь на остановке Ш. Во сколько раз скорость велосипеда меньше скорости автобуса? (Считайте, что автобус и велосипед движутся с постоянными скоростями и останавливаются мгновенно.) }{ 1.25 }{refactoring15.js}
\par{Между остановками Ш. и А. $4,5$ километра прямой трассы. Анастасия выехала на велосипеде из леса между станциями в $225$ метрах от Ш. и увидела, что к Ш. в направлении к А. с постоянной скоростью подъезжает автобус, на который Анастасии нужно успеть. Она заметила, что если она сейчас поедет в сторону Ш., она окажется там одновременно с автобусом. Но и если она поедет в сторону А., она также окажется там одновременно с автобусом, который успеет преодолеть весь участок от Ш. до А., не останавливаясь на остановке Ш. Каково сейчас расстояние (в км) между Анастасией и автобусом? (Считайте, что автобус и велосипед движутся с постоянными скоростями и останавливаются мгновенно.) }{ 0.475 }{refactoring15.js}
\par{Между остановками Х. и А. $3$ километра прямой трассы. Олеся выехала на велосипеде из леса между станциями в $250$ метрах от Х. и увидела, что к Х. в направлении к А. с постоянной скоростью подъезжает автобус, на который Олесе нужно успеть. Она заметила, что если она сейчас поедет в сторону Х., она окажется там одновременно с автобусом. Но и если она поедет в сторону А., она также окажется там одновременно с автобусом, который успеет преодолеть весь участок от Х. до А., не останавливаясь на остановке Х. Каково сейчас расстояние (в км) между Олесей и автобусом? (Считайте, что автобус и велосипед движутся с постоянными скоростями и останавливаются мгновенно.)}{ 0.55 }{refactoring15.js}

\lstinputlisting[caption = refactoring5999.js, label={lst:refactoring5999}, escapechar=|]{code/refactoring5999.js}
\subsubsection*{Примеры генерируемых задач refactoring5999.js}
 
\par{Яхта в $4:00$ вышла из города A в город B, расположенный в $59$ км от A. Пробыв в городе B $5$ часов $7$ минут, яхта отправилась назад и вернулась в город А в $16:00$ того же дня. Определите (в км/ч) скорость течения реки, если известно, что собственная скорость яхты равна $21$ км/ч. }{ 9 }{refactoring5999.js}
\par{Байдарка в $13:00$ вышла из города A в город B, расположенный в $54$ км от A. Пробыв в городе B $3$ часа $18$ минут, байдарка отправилась назад и вернулась в город А в $22:00$ того же дня. Определите (в км/ч) скорость течения реки, если известно, что собственная скорость байдарки равна $19$ км/ч. }{ 1 }{refactoring5999.js}
\par{Баржа в $7:00$ вышла из пункта A в пункт B, расположенный в $22$ км от A. Пробыв в пункте B $3$ часа $8$ минут, баржа отправилась назад и вернулась в пункт А в $16:00$ того же дня. Найдите (в км/ч) скорость течения реки, если известно, что собственная скорость баржи равна $10$ км/ч. }{ 5 }{refactoring5999.js}

\lstinputlisting[caption = refactoring26597.js, label={lst:refactoring26597}, escapechar=|]{code/refactoring26597.js}
\subsubsection*{Примеры генерируемых задач refactoring26597.js}
 
\par{Первый шланг пропускает на $2,75$ литра в минуту меньше, чем второй. Сколько литров воды в минуту пропускает первый шланг, если известно, что бассейн объёмом $567$ литров второй шланг опустошает на $0,66$ минуты быстрее? }{ 47,25 }{refactoring26597.js}
\par{Вторая труба пропускает на $3,5$ литра в минуту больше, чем первая. Сколько литров жидкости в минуту пропускает вторая труба, если бассейн объёмом $147$ литров она заполняет на $4,8$ минуты быстрее, чем первая труба? }{ 12,25 }{refactoring26597.js}
\par{Первый шланг прокачивает на $11,25$ литра в минуту меньше, чем второй. Сколько литров жидкости в минуту прокачивает первый шланг, если известно, что бойлер объёмом $693$ литров он опустошает на $6,3$ минуты дольше, чем второй шланг? }{ 30 }{refactoring26597.js}

\lstinputlisting[caption = refactoring320194.js, label={lst:refactoring320194}, escapechar=|]{code/refactoring320194.js}
\subsubsection*{Примеры генерируемых задач refactoring320194.js}
 
\par{В группе туристов $80$ человек. Их вертолётом в несколько приёмов забрасывают в труднодоступный район по $4$ человека за рейс. Порядок, в котором вертолёт перевозит туристов, случаен. Найдите вероятность того, что турист Ф. полетит семнадцатым рейсом вертолёта. }{ 0,05 }{refactoring320194.js}
\par{В группе туристов $20$ человек. Их вертолётом в несколько приёмов забрасывают в труднодоступный район по $5$ людей за рейс. Порядок, в котором вертолёт перевозит туристов, случаен. Найдите вероятность того, что турист Д. полетит вторым рейсом вертолёта. }{ 0,25 }{refactoring320194.js}
\par{В группе туристов $60$ человек. Их вертолётом в несколько приёмов забрасывают в труднодоступный район по $6$ людей за рейс. Порядок, в котором вертолёт перевозит туристов, случаен. Найдите вероятность того, что турист П. полетит первым рейсом вертолёта. }{ 0,1 }{refactoring320194.js}

\lstinputlisting[caption = refactoring325904.js, label={lst:refactoring325904}, escapechar=|]{code/refactoring325904.js}
\subsubsection*{Примеры генерируемых задач refactoring325904.js}
 
\par{За круглый стол на $11$ стульев в случайном порядке рассаживаются $9$ мальчиков и $2$ девочки. Найдите вероятность того, что обе девочки не будут сидеть рядом. }{ 0,8 }{refactoring325904.js}
\par{За круглый стол на $5$ стульев в случайном порядке рассаживаются $3$ мальчиков и $2$ девочки. Найдите вероятность того, что обе девочки не будут сидеть рядом. }{ 0,5 }{refactoring325904.js}
\par{За круглый стол на $41$ стул в случайном порядке рассаживаются $39$ мальчиков и $2$ девочки. Найдите вероятность того, что обе девочки будут сидеть рядом. }{ 0,05 }{refactoring325904.js}
